\documentclass{article}

% If you need to pass options to natbib, use, e.g.:
% \PassOptionsToPackage{numbers, compress}{natbib}
% before loading neurips_2026.
\usepackage[main]{neurips_2026}

\usepackage[utf8]{inputenc} % allow utf-8 input
\usepackage[T1]{fontenc}    % use 8-bit T1 fonts

\usepackage[american]{babel}
\usepackage{microtype}
\usepackage{mathtools}
\usepackage{amssymb}
\usepackage{amsthm}
\usepackage{algorithm}
\usepackage{algorithmic}
\usepackage{enumitem}
\usepackage{placeins}
\usepackage{float}
\usepackage{booktabs}
\usepackage{graphicx}
\usepackage{xcolor}         % required by NeurIPS checklist macros
\usepackage{hyperref}       % hyperlinks
\usepackage{url}            % simple URL typesetting
\usepackage[capitalize,noabbrev]{cleveref}

% Attempt to make hyperref and algorithmic work together better.
\providecommand{\theHalgorithm}{\arabic{algorithm}}

\theoremstyle{plain}
\newtheorem{theorem}{Theorem}[section]
\newtheorem{proposition}[theorem]{Proposition}
\newtheorem{lemma}[theorem]{Lemma}
\newtheorem{corollary}[theorem]{Corollary}
\theoremstyle{definition}
\newtheorem{definition}[theorem]{Definition}
\newtheorem{assumption}[theorem]{Assumption}
\theoremstyle{remark}
\newtheorem{remark}[theorem]{Remark}

% Double-blind submission: do not include identifying URLs.
\newcommand{\maskimputerepolink}{\texttt{[URL redacted for double-blind review]}}
\newcommand{\ccimputedatasetslink}{\texttt{[URL redacted for double-blind review]}}

% ACM accessibility macro appears in imported figure blocks.
\newcommand{\Description}[1]{}

\setlist[enumerate]{nosep}
\emergencystretch=1em
\newenvironment{compitemize}
  {\begin{itemize}[noitemsep, topsep=0pt, parsep=0pt, partopsep=0pt]}
  {\end{itemize}}

\title{MaskImpute: A Biozero-Aware Masked Autoencoder for Single-Cell RNA-seq Imputation}
\author{Anonymous Author(s)}

\begin{document}
\maketitle
\begin{abstract}
Single-cell RNA sequencing (scRNA-seq) matrices are sparse and noisy, and
observed zeros conflate biological non-expression with technical
under-detection. We propose MaskImpute, a masked denoising autoencoder that uses
different corruption schemes for observed zeros and non-zeros, a count-derived
biozero prior as a soft regularizer, and reconstruction-informed shrinkage on
observed zeros. On a Splatter simulation suite of 13 scenarios with known
pre-dropout ground truth, MaskImpute reduces
mean MSE by 38\% relative to DCA and 35\% relative to scVI, with corresponding
dropout-MSE reductions of 34\% and 32\%, while keeping mean Biozero-MSE below
both neural baselines. The benchmark reports full-matrix, dropout-focused,
biozero-focused, marker-gene, and per-gene normalized reconstruction metrics in
a fixed, dataset-independent configuration.
\end{abstract}

\begin{figure}[H]
\centering
\includegraphics[width=\linewidth]{figures/maskimpute_pipeline.png}
\caption{\textbf{MaskImpute overview.} Starting from logcounts
$\mathbf{X}_{\log}$ and counts $\mathbf{C}$, we estimate a biozero probability
$p_{\mathrm{bio}}(i,j)$ for observed zeros (Section~\ref{sec:biozero}). During
self-supervised training, non-zeros are masked at a fixed rate, while zeros are
masked with probability proportional to $p_{\mathrm{bio}}$; masked zeros are
corrupted by injecting gene-scale noise, and masked non-zeros are replaced by
zero in the scaled space. A fully connected autoencoder reconstructs the scaled
matrix and is trained with the objective
$\mathcal{L}=\mathcal{L}_{\mathrm{mask}}+\lambda_{\mathrm{bio}}\mathcal{L}_{\mathrm{bio}}+\lambda_{\mathrm{nz}}\mathcal{L}_{\mathrm{nz}}$
(Section~\ref{sec:masking}). At inference, we invert the scaling to obtain
$\hat{\mathbf{X}}_{\log}$ and apply reconstruction-informed shrinkage on
observed zeros in the normalized count scale (Section~\ref{sec:zero_shrink}).}
\label{fig:maskimpute_pipeline}
\end{figure}

\section{Introduction}
Single-cell RNA sequencing (scRNA-seq) enables large-scale measurement of RNA expression across individual cells and has become central to studying cellular heterogeneity, developmental trajectories, disease states, and perturbation responses \citep{Luecken2019BestPractices, Hicks2018MissingData, VanDeSande2023DrugDiscovery}. The resulting gene-by-cell matrix is sparse and noisy because only a small fraction of molecules are captured and sequenced, and additional technical factors (e.g., batch effects) perturb measurements \citep{Brennecke2013TechnicalNoise, Hicks2018MissingData}. A central difficulty is the interpretation of zeros: an observed zero is a valid measurement that can arise from true biological non-expression or from limited molecular sampling (often termed ``dropout''), and disentangling these sources from expression alone is generally not identifiable \citep{SarkarStephens2021SeparatingMeasurement, Svensson2020NotZeroInflated, Jiang2022ZeroInflationReview, Linderman2022ZeroPreserving}. This uncertainty complicates downstream analyses such as differential expression, trajectory inference, and cell-state interpretation \citep{Hou2020ImputationBenchmark, Luecken2019BestPractices}.

A large literature proposes denoising and imputation methods that treat scRNA-seq as a structured inference problem: borrow information across related cells and genes to reduce measurement noise. Graph diffusion methods such as MAGIC smooth expression by propagating information over a cell--cell neighborhood graph \citep{vanDijk2018MAGIC}; they can reveal global structure but may over-smooth within-state variability and inflate gene--gene correlations when diffusion is aggressive \citep{Huang2018SAVER, Hou2020ImputationBenchmark}. Probabilistic approaches such as SAVER and scImpute model aspects of the measurement process and/or dropout explicitly \citep{Huang2018SAVER, Li2018scImpute}, but their performance depends on modeling assumptions and on how well ``dropout'' can be operationalized for a given protocol \citep{Svensson2020NotZeroInflated, Jiang2022ZeroInflationReview}. Deep autoencoder methods (e.g., DCA and AutoClass) scale to large datasets and learn nonlinear manifolds \citep{Eraslan2019DCA, Li2022AutoClass}, yet reconstructions are typically point estimates and can still blur rare states or distort variance if treated as ground truth \citep{Hou2020ImputationBenchmark, Linderman2022ZeroPreserving}.

We introduce \textbf{MaskImpute}, a masked denoising autoencoder designed around
an uncertainty-aware view of zeros. Instead of treating every zero as a missing
value to be filled, MaskImpute combines (i) a masking strategy with different
corruption schemes for observed zeros vs.\ observed non-zeros, (ii) a
count-derived biozero prior used only as a soft regularizer, and (iii) a
reconstruction-informed shrinkage step applied only on observed zeros to reduce
over-imputation while retaining recovery of likely technical under-detection.

This work makes the following contributions:
\begin{compitemize}
  \item We adapt masked self-supervised learning to scRNA-seq denoising with different masking/corruption schemes for zeros and non-zeros.
  \item We introduce a biozero-aware regularization and masking scheme that encourages zero-preserving behavior while still allowing recovery of likely technical under-detection.
  \item We introduce a reconstruction-informed post-hoc shrinkage mechanism that modulates imputed values only at observed zeros based on the model's reconstructed signal.
  \item We benchmark MaskImpute against representative baselines on a synthetic suite with known ground truth, reporting full-matrix, dropout-focused, biozero-focused, marker-gene, and per-gene normalized errors.
  \item We provide a reproducible evaluation protocol that keeps MaskImpute settings fixed across synthetic scenarios and separates imputation accuracy from downstream task-specific post-processing.
\end{compitemize}

\section{Related Works}

\subsection{Statistical and Probabilistic Models}
scRNA-seq denoising sits at the intersection of representation learning and uncertainty-aware inference: methods attempt to estimate latent expression under an imperfect measurement process, while avoiding the creation of spurious signal that can mislead downstream reasoning \citep{SarkarStephens2021SeparatingMeasurement, Hou2020ImputationBenchmark, Linderman2022ZeroPreserving}.

Early methods focused on explicit probabilistic models of counts. SAVER \citep{Huang2018SAVER} assumes a Poisson measurement model with a Gamma prior (negative binomial marginal) and uses empirical Bayes regression to estimate gene- and cell-specific priors, producing posterior means and uncertainty estimates. scImpute \citep{Li2018scImpute} uses mixture modeling to estimate dropout probabilities and selectively imputes only entries deemed likely to be technical, borrowing information from similar cells. These approaches can be conservative and interpretable, but they depend on distributional assumptions and on how well the model captures protocol-specific noise (including the debated role of zero inflation) \citep{Svensson2020NotZeroInflated, Jiang2022ZeroInflationReview}.

\subsection{Graph-Based and Diffusion Methods}
Graph-based methods assume that biologically similar cells lie on (or near) a
low-dimensional manifold and that neighborhoods can be used to borrow signal
across cells. MAGIC \citep{vanDijk2018MAGIC} constructs a $k$-nearest-neighbor
affinity graph (often in a PCA space), normalizes it into a Markov transition
matrix $M$, and imputes by diffusion,
$\mathbf{X}' = M^{t}\mathbf{X}$, where $t$ is the diffusion time. Such methods
can be effective for visualization and smooth trajectories, but the global
averaging implicit in diffusion can blur rare states and inflate correlations
if neighborhood or diffusion parameters are too aggressive
\citep{Huang2018SAVER, Hou2020ImputationBenchmark, Linderman2022ZeroPreserving}.

ccImpute \citep{Malec2022ccImpute} offers an alternative
that builds a consensus similarity matrix from repeated K-means clustering
(inspired by SC3 \citep{Kiselev2017SC3}) and imputes likely dropouts by
consensus-weighted averaging over similar cells. Compared with diffusion,
consensus approaches can reduce sensitivity to a single neighborhood choice,
but they still inherit clustering biases and can sharpen boundaries; thus
improvements in clustering metrics after imputation should be interpreted
cautiously \citep{Hou2020ImputationBenchmark}.

\subsection{Deep Learning and Autoencoders}
Autoencoder architectures scale to large datasets and can model nonlinear
structure, making them attractive for scRNA-seq denoising. DCA
\citep{Eraslan2019DCA} replaces a squared-error reconstruction loss with a
count likelihood (NB/ZINB) so that the decoder outputs distribution parameters
(e.g., mean and dispersion, and optionally zero-inflation). This can improve
robustness to heavy-tailed count noise, but the interpretation of zero-inflation
and per-entry dropout probabilities is limited by identifiability and by the
ongoing debate over whether additional zero-inflation is needed for common UMI
protocols \citep{Svensson2020NotZeroInflated, Jiang2022ZeroInflationReview}.
More broadly, neural denoisers typically output point estimates that can be
over-interpreted as ``ground truth'' and can distort variance or downstream
statistics if used uncritically \citep{Hou2020ImputationBenchmark, Linderman2022ZeroPreserving}.

scVI \citep{Lopez2018scVI} is a variational autoencoder that models scRNA-seq
counts with an explicit generative process and has become a common baseline for
scalable representation learning and denoising. In contrast to neural
approaches, ALRA \citep{Linderman2022ALRA} performs an adaptively-thresholded
low-rank approximation that explicitly preserves zeros, providing a strong
zero-preserving imputation baseline.

AutoClass \citep{Li2022AutoClass} augments an autoencoder with a classifier
branch trained on pseudo-labels from pre-clustering, encouraging the bottleneck
to preserve population structure. This can improve separability, but it
introduces dependence on the pre-clustering signal and can propagate
confounders (e.g., batch structure) present in pseudo-labels
\citep{Hou2020ImputationBenchmark}. MaskImpute adopts a self-supervised masking
objective to learn structure without requiring explicit pseudo-labels.

\subsection{Self-Supervised and Masked Learning}
Masked self-supervision has achieved strong results in vision and language by
masking parts of the input and training models to reconstruct them, which
encourages learning global dependencies rather than memorizing local statistics
\citep{Devlin2019BERT, He2022MAE}. Recent work has also applied masked
autoencoders to scRNA-seq for representation learning
\citep{Fang2024scMAE}. However, in scRNA-seq the distinction between an
observed zero and a missing value is itself uncertain, and naive reconstruction
objectives can encourage over-imputation of genuine biological zeros
\citep{SarkarStephens2021SeparatingMeasurement, Linderman2022ZeroPreserving}.
MaskImpute adapts masking to this setting by (i) using different masking rates
for zeros and non-zeros, and (ii) incorporating a biozero-aware prior and
regularization so that reconstruction pressure is reduced on zeros likely to be
biological non-expression.

\section{Methods}
\label{sec:methods}
\subsection{Synthetic scRNA-seq Datasets}
\label{sec:synthetic_data}
We benchmarked MaskImpute against other approaches on deterministic synthetic
scRNA-seq datasets generated with the Splat simulator from the
\texttt{Splatter} package \citep{zappia2017splatter}. The goal is controlled
evaluation in a setting where both the observed counts and pre-dropout ground
truth are available.

\paragraph{Summary.}
We simulate 13 scenarios with $G=1000$ genes spanning grouped populations, batch
effects, and trajectory topologies. Each scenario is generated deterministically
with fixed seeds and split into a tuning split ($N=1{,}000$ cells), a held-out
test split ($N=5{,}000$ cells), and a runtime-scaling split
($N\in\{10{,}000,\ldots,100{,}000\}$) on a representative dropout scenario.
Datasets store observed counts and pre-dropout ground truth, along with
log2(1+CP10k) assays; marker genes are those with differential-expression
factors not equal to 1 in any group. Full simulator parameters and the scenario
catalog are provided in Appendix~\ref{app:synthetic_details}.

\paragraph{Notation.}
In this work, $\mathbf{X}_{\log}\in\mathbb{R}_{\ge 0}^{N\times G}$ denotes the
observed log-normalized matrix, $\mathbf{X}^{\star}_{\log}$ the corresponding
synthetic ground truth (when available), and $\mathbf{X}'_{\log}$ the
reconstructed matrix produced by MaskImpute. We write $x_{ij}$, $x^{\star}_{ij}$,
and $x'_{ij}$ for their respective entries. We denote the (unnormalized) count
matrix by $\mathbf{C}=\{c_{ij}\}$.

\subsection{Model Architecture}
MaskImpute is a masked denoising autoencoder trained with a biozero-aware
corruption and regularization scheme to reduce over-imputation at observed
zeros. Let $\tilde{\mathbf{s}}_i\in\mathbb{R}^G$ denote the corrupted, scaled
input for cell $i$ (Sections~\ref{sec:scaling}--\ref{sec:masking}). The encoder
$f_\theta:\mathbb{R}^{G}\!\to\!\mathbb{R}^{d_{\mathrm{lat}}}$ maps $\tilde{\mathbf{s}}_i$ to a
latent representation $\mathbf{h}_i=f_\theta(\tilde{\mathbf{s}}_i)$, and a
decoder $g_\phi:\mathbb{R}^{d_{\mathrm{lat}}}\!\to\!\mathbb{R}^{G}$ produces the reconstruction
  \[
\hat{\mathbf{s}}_i=g_\phi(\mathbf{h}_i)=g_\phi(f_\theta(\tilde{\mathbf{s}}_i)).
  \]

  We use a symmetric fully connected architecture with encoder widths
  $128 \rightarrow 64 \rightarrow d_{\mathrm{lat}}$ and decoder widths
  $d_{\mathrm{lat}} \rightarrow 64 \rightarrow 128 \rightarrow G$, with bottleneck dimension
  $d_{\mathrm{lat}}{=}16$. Each hidden block applies a linear transformation, layer
  normalization, and a SiLU nonlinearity. This yields a compact latent
  representation while preserving enough capacity to model nonlinear
  gene--gene structure.

  Let $\mathbf{S}=\mathrm{scale}(\mathbf{X}_{\log})$ and $\hat{\mathbf{S}}$
  denote the scaled inputs and reconstructions stacked across cells.
  Reconstructions in log space are obtained by inverse scaling,
  \[
  \mathbf{X}'_{\log}=\mathrm{invscale}(\hat{\mathbf{S}}).
  \]
  All entries, including observed non-zeros, are reconstructed by the same model.
  This makes the method a full-matrix imputer rather than a zero-only editing
  procedure.

\subsection{Gene Scaling}
\label{sec:scaling}
We apply a robust, monotone per-gene scaling with three steps:
\[
\begin{aligned}
\text{(Winsorize)}\quad &\bar{x}_{ij}=\min(\max(x_{ij},q^{(2)}_j),q^{(99.5)}_j),\\
\text{(Standardize)}\quad &z_{ij}=\frac{\bar{x}_{ij}-\mu_j}{\sigma_j},\qquad
\sigma_j\leftarrow\max(\sigma_j,\varepsilon),\\
\text{(Rescale)}\quad &s_{ij}=2\cdot\frac{z_{ij}-z^{\min}_j}{z^{\max}_j-z^{\min}_j}-1,\\
&\qquad z^{\max}_j-z^{\min}_j\ge\varepsilon.
\end{aligned}
\]
We denote the scaled matrix by $\mathbf{S}$ and use
$\mathrm{scale}(\cdot)$ / $\mathrm{invscale}(\cdot)$ for forward/inverse
transforms.
Here $\varepsilon>0$ is an arbitrarily small constant used for numerical
stability (we use $\varepsilon=10^{-8}$).
We define $z^{\min}_j=\min_i z_{ij}$ and $z^{\max}_j=\max_i z_{ij}$.

\subsection{Masked Denoising Objective}
\label{sec:masking}
  For each mini-batch $\mathcal{B}$ of cells, we sample corruption masks with
  different rates for observed zeros vs.\ observed non-zeros. Let $c_{ij}$ denote
  the observed count and let $x_{ij}$ denote the corresponding log-normalized
  value. Define the observed-zero set $\mathcal{Z}=\{(i,j):c_{ij}=0\}$ and the
  observed non-zero set $\mathcal{N}=\{(i,j):c_{ij}>0\}$. We sample two Bernoulli
  masks:
  \[
  \begin{aligned}
  m^{(0)}_{ij} &\sim \mathrm{Bernoulli}(\pi^{(0)}_{ij}),\qquad
  \pi^{(0)}_{ij}=p_{\mathrm{zero}}\,p_{\mathrm{bio}}(i,j)
  &&\text{for }(i,j)\in\mathcal{Z},\\
  m^{(\mathrm{nz})}_{ij} &\sim \mathrm{Bernoulli}(p_{\mathrm{nz}})
  &&\text{for }(i,j)\in\mathcal{N},
  \end{aligned}
  \]
  where $p_{\mathrm{bio}}(i,j)$ is the biozero probability for an observed zero
  (Section~\ref{sec:biozero}). We use $(p_{\mathrm{zero}},p_{\mathrm{nz}})=(0.01,0.05)$.

  Masked non-zeros are replaced by zero in the scaled space. Masked zeros are
  replaced by count-derived noise as described below. Unmasked entries are left
  unchanged.

  \paragraph{Noise for masked zeros.}
  For masked zero entries, we inject noise in log space using the observed count
  range per gene. Let $c^{\max}_j=\max_i c_{ij}$. For each masked zero entry we
  sample
  \[
  \eta_{ij}\sim\mathrm{Uniform}(0,\nu),\quad \nu=0.02,\qquad
  \tilde{x}_{ij}=\log_2(1+\eta_{ij}c^{\max}_j),
  \]
  then apply $\mathrm{scale}(\cdot)$ to obtain the corrupted scaled input value.

\paragraph{Loss.}
  Let $s_{ij}=[\mathrm{scale}(\mathbf{X}_{\log})]_{ij}$ and
  $\hat{s}_{ij}$ denote the reconstruction in the scaled space, and let
  $r_{ij}=\hat{s}_{ij}-s_{ij}$ be the residual. We optimize a weighted masked MSE:
  \[
  \mathcal{L}_{\mathrm{mask}}=
  \frac{\omega_{\mathrm{bio}}\sum_{(i,j)\in\mathcal{Z}} m^{(0)}_{ij}\,r_{ij}^2
        +\omega_{\mathrm{nz}}\sum_{(i,j)\in\mathcal{N}} m^{(\mathrm{nz})}_{ij}\,r_{ij}^2}
       {\omega_{\mathrm{bio}}\sum_{(i,j)\in\mathcal{Z}} m^{(0)}_{ij}
        +\omega_{\mathrm{nz}}\sum_{(i,j)\in\mathcal{N}} m^{(\mathrm{nz})}_{ij}},
  \]
  with weights $(\omega_{\mathrm{bio}},\omega_{\mathrm{nz}})=(1.0,2.0)$.

  We add two auxiliary reconstruction terms. The biozero regularizer encourages
  predicted values at likely biological zeros to be close to the gene-wise scaled
  zero value:
  \[
  \begin{aligned}
  \mathcal{L}_{\mathrm{bio}}
  &=\frac{\sum_{(i,j)\in\mathcal{Z}} p_{\mathrm{bio}}(i,j)\,(\hat{s}_{ij}-s_{0j})^2}
         {\sum_{(i,j)\in\mathcal{Z}} p_{\mathrm{bio}}(i,j)},\qquad
  s_{0j}=[\mathrm{scale}(\mathbf{0})]_j.
  \end{aligned}
  \]
  Because the model reconstructs the full matrix (including observed non-zeros),
  we also include an observed-expression reconstruction penalty over $\mathcal{N}$,
  \[
  \mathcal{L}_{\mathrm{nz}}=
  \frac{\sum_{(i,j)\in\mathcal{N}} r_{ij}^2}{|\mathcal{N}|}.
  \]
  The final objective is
  \[
  \mathcal{L}=\mathcal{L}_{\mathrm{mask}}+\lambda_{\mathrm{bio}}\,\mathcal{L}_{\mathrm{bio}}
  +\lambda_{\mathrm{nz}}\,\mathcal{L}_{\mathrm{nz}},
  \]
  with $(\lambda_{\mathrm{bio}},\lambda_{\mathrm{nz}})=(1.2,0.75)$.

\subsection{Biological Zero Probability Estimation}
  \label{sec:biozero}
  We estimate a prior probability $p_{\mathrm{bio}}(i,j)$ that an observed zero
  corresponds to true biological non-expression. We use a negative-binomial
  model for latent true counts with an explicit dropout mechanism. Let $t_{ij}$
  denote a latent (pre-dropout) count and $d_{ij}$ a dropout indicator:
  \[
  t_{ij} \sim \mathrm{NB}(\mu_{ij},\phi_j),\qquad
  d_{ij} \sim \mathrm{Bernoulli}(\delta_{ij}),\qquad
  c_{ij} = (1-d_{ij})\,t_{ij}.
  \]
  Let $f_{ij}(0)=\Pr(t_{ij}=0\mid \mu_{ij},\phi_j)$ denote the NB zero probability.
  Then for an observed zero (noting that $c_{ij}=0$ whenever $t_{ij}=0$,
  regardless of dropout),
  \[
  p_{\mathrm{bio}}(i,j)=\Pr(t_{ij}=0\mid c_{ij}=0)
  =\frac{f_{ij}(0)}{\delta_{ij}+(1-\delta_{ij})f_{ij}(0)}.
  \]
  We set $p_{\mathrm{bio}}(i,j)=0$ when $c_{ij}>0$.

  We model $\mu_{ij}$ using a gene mean $\mu_j$ and a library-size factor,
  \[
  \mu_{ij}=\ell_i\mu_j,\qquad
  \ell_i=\frac{\sum_j c_{ij}}{\operatorname*{median}_k\left(\sum_j c_{kj}\right)}.
  \]
  Dropout probabilities $\delta_{ij}$ are modeled as a smooth decreasing function
  of $\log \mu_{ij}$ using a fitted logistic curve to empirical zero rates, and
  gene-wise dispersion parameters $\phi_j$ are estimated using a shrinkage
  moment-based estimator on observed counts. This prior is used strictly as a
  soft regularizer and as a modulation signal for zero masking; it is not assumed
  to identify ground-truth biological zeros in real data.


\subsection{Training Details}
We train MaskImpute in a transductive self-supervised setting: the model is fit
on the full observed matrix and learns to reconstruct randomly corrupted
entries. Unless otherwise specified, we optimize the objective in
Section~\ref{sec:masking} with Adam (learning rate $10^{-4}$, weight decay
$10^{-4}$), batch size 32, for 400 epochs on GPU. Random seeds are fixed for
reproducibility. Biozero probabilities $p_{\mathrm{bio}}(i,j)$ are computed once
from observed counts (Section~\ref{sec:biozero}) and held fixed during training.
All hyperparameters are held constant across datasets and splits.

\subsection{Reconstruction-Informed Zero Shrinkage}
\label{sec:zero_shrink}
The autoencoder outputs a full reconstruction $\hat{\mathbf{X}}_{\log}$ (all
entries, including observed non-zeros). To reduce over-imputation at observed
zeros, we apply a lightweight post-hoc shrinkage \emph{only} on entries with
$c_{ij}=0$. Let $\hat{u}_{ij}=2^{\hat{x}_{ij}}-1$ denote the reconstructed
CP10k-normalized count scale corresponding to $\hat{x}_{ij}$. For each gene $j$,
let $\tilde{u}_j$ be a robust reference level, taken as the median of
$\hat{u}_{ij}$ over cells where the gene is observed as expressed ($c_{ij}>0$),
with a fallback to the 75th percentile of $\{\hat{u}_{ij}\}_i$ when needed.
We define a reconstruction-derived biozero weight
\[
w^{(s)}_{ij}=\frac{1}{1+\hat{u}_{ij}/\max(\tilde{u}_j,\varepsilon)}\qquad\text{for }c_{ij}=0,
\]
and set $w^{(s)}_{ij}=0$ for $c_{ij}>0$. The final imputed normalized counts are
\[
\hat{u}^{(\mathrm{shrink})}_{ij}=\hat{u}_{ij}\big(1-\alpha w^{(s)}_{ij}\big)\qquad\text{for }c_{ij}=0,
\]
with shrinkage strength $\alpha=0.3$, while entries with $c_{ij}>0$ are left
unchanged by this step. We report final reconstructions in log space as
$x'_{ij}=\log_2(1+\hat{u}^{(\mathrm{shrink})}_{ij})$.

\subsection{Evaluation Metrics}
\label{sec:evaluation}
\paragraph{(i) Reconstruction error (synthetic test split).}
On the synthetic test split ($N=5000$ cells per scenario), we compare
$\mathbf{X}'_{\log}$ to the ground truth $\mathbf{X}^{\star}_{\log}$ (both in
log2(1+CP10k) space). We report mean squared error (MSE) and mean absolute error
(MAE),
\[
\mathrm{MSE}=\frac{1}{NG}\sum_{i,j}\big(x^{\star}_{ij}-x'_{ij}\big)^2,\qquad
\mathrm{MAE}=\frac{1}{NG}\sum_{i,j}\big|x^{\star}_{ij}-x'_{ij}\big|.
\]
We also report subset metrics on (a) dropout-like entries
$D=\{(i,j):x_{ij}=0\wedge x^{\star}_{ij}>0\}$, (b) biozero entries
$Z_{\mathrm{bio}}=\{(i,j):x^{\star}_{ij}=0\}$, and (c) marker-gene entries
$C=\{(i,j):j\in\mathcal{J}_M\}$, where $\mathcal{J}_M$ is defined from
\texttt{DEFacGroup*} (Section~\ref{sec:synthetic_data}). For each subset, the
metric is the corresponding average over the subset; if a subset is empty we
report it as NA.
For methods that output imputed counts rather than logcounts, we normalize the
imputed count matrix to the dataset \texttt{targetSum} using its \emph{imputed}
per-cell library sizes, and then apply the $\log_2(1+x)$ transform; we do not
use the ground-truth library sizes (\texttt{libSizeTrue}). Unless otherwise
specified, stochastic methods are run for five random seeds and results are
reported as mean $\pm$ standard deviation.

\paragraph{(ii) Per-gene normalized RMSE (gNRMSE).}
To reduce sensitivity to the scale of highly expressed genes, we report
per-gene normalized RMSE,
\[
\mathrm{gNRMSE}=\frac{1}{G}\sum_{j=1}^G\frac{\sqrt{\frac{1}{N}\sum_i\left(x^{\star}_{ij}-x'_{ij}\right)^2}}
{\max\left(\mathrm{sd}(\{x^{\star}_{ij}\}_{i=1}^N),\varepsilon\right)},
\]
and its marker-gene variant, averaged over $j\in\mathcal{J}_M$.

\paragraph{(iii) Runtime scaling.}
To evaluate scalability, we measure wall-clock runtime on the scaling split for
\[
N\in\{10000,20000,\ldots,100000\}
\]
on the representative scenario \texttt{groups\_balanced\_moderate\_drop}. We
report mean runtime (and standard deviation across datasets when applicable).
All runtime benchmarks and reported runtimes were obtained on an NVIDIA DGX
H100 system and all methods were benchmarked on the same hardware
(Appendix~\ref{app:runtime}).

\paragraph{(iv) Real-data label agreement.}
For labeled real datasets, no pre-dropout ground truth is available. We therefore
evaluate only whether reconstructed matrices preserve externally annotated cell
types. We use the full Baron human pancreas dataset \citep{Baron2016Pancreas},
the full Zeisel mouse cortex dataset \citep{Zeisel2015Cortex}, and the 10x
PBMC68k dataset with the published barcode-level annotations from the original
analysis \citep{Zheng2017TenX}. To avoid supervised marker selection, we select
the top 1,000 highly variable genes from the observed matrix only, using
mean-binned dispersion of log-normalized expression. Labels are used only after
imputation for evaluation.
For computationally tractable repeated evaluation, PBMC68k is evaluated on a
fixed random 12,000-cell subset. For each imputed matrix, we standardize genes,
compute PCA embeddings with $d\in\{20,50,100\}$ components, build a shared
nearest-neighbor graph with $k\in\{10,15,30\}$, and run Leiden clustering
\citep{Traag2019Leiden} at resolutions
$\{0.2,0.5,1.0,1.5,2.0\}$. We repeat the procedure for five seeds and report
the mean and standard deviation across this grid. We also report a
cluster-count-matched summary that selects, for each seed, the grid point whose
number of clusters is closest to the number of annotated cell types; this uses
only the number of labels, not label identities. When batch-like metadata are
available (Baron donor, Zeisel cell-id prefix, PBMC barcode suffix), we report
cluster--batch ARI as a diagnostic, where lower values indicate weaker
alignment between clusters and batch. DCA is run with a 300-second wall-clock
budget per dataset; timeouts are recorded explicitly rather than omitted.
This evaluation remains a downstream sanity check rather than a direct
denoising metric, and it can favor conservative methods when raw expression
already separates annotated populations.

\paragraph{(v) Synthetic downstream label agreement and calibration.}
For the phase-6 diagnostic analyses, we also evaluate clustering on the
synthetic test split using the simulator-provided group or path labels. Labels
are used only for evaluation. We use the same PCA--SNN--Leiden procedure as in
the real-data analysis, but with a smaller robustness grid
$d\in\{20,50\}$, $k\in\{10,15,30\}$, and resolutions
$\{0.2,0.5,1.0\}$. We report both grid-averaged ARI/NMI and a
cluster-count-matched summary. We additionally examine a one-parameter,
label-free residual calibration
\[
\mathbf{X}^{(\lambda)}_{\log}=(1-\lambda)\mathbf{X}'_{\log}
+\lambda\mathbf{X}_{\log},
\]
where $\mathbf{X}_{\log}$ is the observed log-normalized matrix. This transform
does not use labels or dataset-specific fitting and is evaluated only as a
simple calibration of the final output.

\subsection{Algorithm Outline}
Figure~\ref{fig:maskimpute_pipeline} provides a high-level view of the MaskImpute
pipeline, from biozero probability estimation to masked corruption and
reconstruction. Algorithm~\ref{alg:maskimpute} gives the complete
training and imputation procedure, including mask sampling, corruption, and the
loss minimized by MaskImpute.

\begin{algorithm}[H]
\caption{MaskImpute training and imputation.}
\label{alg:maskimpute}
\begin{algorithmic}[1]
\raggedright
\STATE \textbf{Input:} logcounts $\mathbf{X}_{\log}$, counts $\mathbf{C}$,
scaler $\mathrm{scale}/\mathrm{invscale}$, masking rates
$p_{\mathrm{zero}},p_{\mathrm{nz}}$, noise level $\nu$, shrinkage strength
$\alpha$, and epochs $E$.
\STATE Compute biozero probabilities $p_{\mathrm{bio}}(i,j)$ from counts via
Section~\ref{sec:biozero}.
\STATE Scale $\mathbf{S}\leftarrow\mathrm{scale}(\mathbf{X}_{\log})$.
\FOR{epoch $=1$ to $E$}
\FOR{mini-batch $\mathcal{B}$}
\STATE Sample masks on entries: $m^{(0)}_{ij}\sim\mathrm{Bernoulli}(p_{\mathrm{zero}}p_{\mathrm{bio}}(i,j))$
for observed zeros ($c_{ij}=0$), and $m^{(\mathrm{nz})}_{ij}\sim\mathrm{Bernoulli}(p_{\mathrm{nz}})$
for observed non-zeros ($c_{ij}>0$).
\STATE Corrupt inputs: masked non-zeros $\tilde{s}_{ij}\leftarrow 0$; masked
zeros $\tilde{x}_{ij}=\log_2(1+\eta_{ij}c^{\max}_j)$ with
$\eta_{ij}\sim\mathrm{Uniform}(0,\nu)$ and $c^{\max}_j=\max_i c_{ij}$, then
$\tilde{s}_{ij}\leftarrow\mathrm{scale}(\tilde{x}_{ij})$.
\STATE Forward pass: $\hat{\mathbf{s}}_i=g_\phi(f_\theta(\tilde{\mathbf{s}}_i))$.
\STATE Update parameters by minimizing $\mathcal{L}$ (Section~\ref{sec:masking}).
\ENDFOR
\ENDFOR
\STATE Reconstruct $\hat{\mathbf{X}}_{\log}\leftarrow\mathrm{invscale}(\hat{\mathbf{S}})$.
\STATE Apply reconstruction-informed shrinkage on observed zeros in the
normalized count scale (Section~\ref{sec:zero_shrink}) to obtain
$\mathbf{X}'_{\log}$.
\STATE \textbf{Output:} reconstructed logcounts $\mathbf{X}'_{\log}$.
\end{algorithmic}
\end{algorithm}
\FloatBarrier

\section{Experiments and Results}
\label{sec:experiments}

\subsection{Denoising Accuracy on Synthetic Data}
\label{sec:denoising_results}
We evaluate imputation accuracy on the 13 synthetic scenarios described in
Section~\ref{sec:synthetic_data} on the held-out test split ($N=5{,}000$ cells
per scenario). We compare reconstructed logcounts $\mathbf{X}'_{\log}$ against
the synthetic ground truth $\mathbf{X}^{\star}_{\log}$ using the metrics in
Section~\ref{sec:evaluation}. In addition to overall reconstruction error, we
report subset errors on dropout-like entries (observed zeros with non-zero
ground truth), biozero entries (ground-truth zeros), and marker-gene entries.
Unless otherwise stated, we run each method for 5 repeats with different random
seeds per scenario. Baselines with available generated results include DCA,
scVI, ALRA, MAGIC, AutoClass, ccImpute, SAVER, and the un-denoised observed
logcounts (Baseline).
scVI is run with a negative-binomial likelihood, two hidden layers, latent
dimension 10, and 400 training epochs; its normalized expression output is
scaled to the benchmark CP10k target before log transformation. ALRA is run as a
fixed adaptive low-rank reconstruction baseline on log-normalized expression,
with rank selected from the singular-value spectrum and gene-wise thresholding
of low reconstructed values.

\begin{table*}[t]
\caption{Synthetic test split benchmark summary across 13 scenarios and 5 seeds.
Entries report means across scenario-level estimates; $\pm$ values are standard
deviations across scenarios. Each scenario-level estimate averages 5 random
seeds. Baseline denotes the un-denoised observed logcounts (identity). Lower is
better.}
\label{tab:benchmark_core}
\centering
\footnotesize
\begin{tabular}{@{}lrrrrr@{}}
\toprule
Method & MSE & Dropout-MSE & Biozero-MSE & MAE & gNRMSE \\
\midrule
MaskImpute & \textbf{0.2242 $\pm$ 0.0399} & \textbf{0.3053 $\pm$ 0.0374} & 0.0133 $\pm$ 0.0055 & \textbf{0.3485 $\pm$ 0.0320} & \textbf{1.1592 $\pm$ 0.0870} \\
DCA & 0.3598 $\pm$ 0.1306 & 0.4652 $\pm$ 0.1373 & 0.0164 $\pm$ 0.0243 & 0.4504 $\pm$ 0.0879 & 1.3867 $\pm$ 0.2161 \\
scVI & 0.3426 $\pm$ 0.1194 & 0.4511 $\pm$ 0.1292 & 0.0179 $\pm$ 0.0225 & 0.4355 $\pm$ 0.0789 & 1.3502 $\pm$ 0.1948 \\
ALRA & 0.6917 $\pm$ 0.2917 & 0.7912 $\pm$ 0.2646 & 0.0103 $\pm$ 0.0205 & 0.6498 $\pm$ 0.1566 & 1.8498 $\pm$ 0.3737 \\
MAGIC & 0.6411 $\pm$ 0.2850 & 0.7796 $\pm$ 0.2693 & 0.0104 $\pm$ 0.0195 & 0.6256 $\pm$ 0.1557 & 1.7747 $\pm$ 0.3693 \\
AutoClass & 0.7281 $\pm$ 0.3038 & 0.8621 $\pm$ 0.2971 & 0.0098 $\pm$ 0.0204 & 0.6721 $\pm$ 0.1593 & 1.9457 $\pm$ 0.3888 \\
ccImpute & 0.6983 $\pm$ 0.4751 & 1.4643 $\pm$ 0.6006 & 0.0095 $\pm$ 0.0308 & 0.5068 $\pm$ 0.2240 & 1.6406 $\pm$ 0.6393 \\
SAVER & 1.2208 $\pm$ 0.5005 & 2.6786 $\pm$ 0.3557 & 0.0020 $\pm$ 0.0064 & 0.6745 $\pm$ 0.2190 & 2.4384 $\pm$ 0.7013 \\
Baseline & 1.5236 $\pm$ 0.5768 & 3.3665 $\pm$ 0.3281 & \textbf{0.0000 $\pm$ 0.0000} & 0.7343 $\pm$ 0.2569 & 2.7044 $\pm$ 0.8656 \\
\bottomrule
\end{tabular}

\end{table*}

Table~\ref{tab:benchmark_core} summarizes average performance across scenarios.
Among the baselines with available results, MaskImpute achieves the lowest mean
MSE and dropout-MSE. Relative to DCA, MaskImpute reduces overall MSE by 38\% and
dropout-MSE by 34\% in this benchmark; relative to scVI, it reduces overall MSE
by 35\% and dropout-MSE by 32\%. Marker-gene subset errors are reported in Appendix~\ref{app:per_scenario}
(Table~\ref{tab:benchmark_marker}).
MaskImpute also has lower mean Biozero-MSE than DCA (0.0133 vs. 0.0164) and
scVI (0.0133 vs. 0.0179), but the DCA comparison is not uniform across
scenarios. Biozero-focused errors are
lowest for conservative strategies that preserve observed zeros (Baseline,
SAVER), which is expected because biozero entries have ground-truth zero
expression. MaskImpute therefore prioritizes dropout recovery and full-matrix
accuracy while maintaining Biozero-MSE below the neural baselines.
We report per-scenario MSE and marker-gene subset errors in
Appendix~\ref{app:per_scenario}, along with scenario-level wins and paired
Wilcoxon signed-rank tests against DCA.

We note that gNRMSE values can exceed 1 because the metric normalizes by each
gene's standard deviation, including low-variance genes where even small
absolute errors translate into large relative error. In our benchmark, MaskImpute
still improves gNRMSE relative to the baselines with available results
(Table~\ref{tab:benchmark_core}).
Appendix~\ref{app:per_scenario} (Figure~\ref{fig:mse_5000_bars}) summarizes
benchmark metrics aggregated across scenarios (mean $\pm$ standard deviation).

\subsection{Ablations, Calibration, and Real-Data Checks}
\label{sec:phase6_results}
Appendix~\ref{app:ablations} reports component ablations, sensitivity analyses,
and a label-free residual-calibration sweep. The ablations show that the
biozero regularizer is the primary component protecting biological zeros:
removing it reduces aggregate MSE but increases Biozero-MSE from 0.013 to
0.095. Removing zero shrinkage has a smaller but consistent effect, increasing
Biozero-MSE to 0.018. A simple residual calibration
$X^{(\lambda)}=(1-\lambda)X_{\mathrm{MI}}+\lambda X_{\mathrm{obs}}$ can improve
the diagnostic efficacy score; $\lambda=0.05$ is the safest synthetic-error
trade-off, while $\lambda=0.10$ gives the strongest synthetic downstream
ARI/NMI in the screened grid.

Appendix~\ref{app:real_data} gives downstream label-agreement results on Baron
human pancreas, Zeisel mouse cortex, and PBMC68k. These real-data experiments do
not provide pre-dropout ground truth and are therefore reported as
structure-preservation checks rather than direct imputation-accuracy evidence.
MaskImpute has the best cluster-count-matched Baron ARI/NMI among evaluated
methods and the best matched Zeisel NMI among imputation methods, but PBMC68k
does not show a MaskImpute clustering advantage. The synthetic ground-truth
benchmark remains the primary evidence for imputation accuracy.

\section{Limitations}
MaskImpute introduces additional hyperparameters (masking rates, regularization,
shrinkage strength) whose interactions may be dataset-dependent; the sensitivity
analysis in Section~\ref{sec:phase6_results} is diagnostic and should be
expanded to additional datasets before final deployment. The count-derived
biozero prior is used only as a soft regularizer and masking modulator, and may
be misspecified for some
protocols or preprocessing pipelines. The real-data experiments use annotated
cell-type labels as a downstream proxy, not ground-truth molecular
reconstructions, and therefore can reward conservative behavior on marker-rich
datasets. Finally, the current runtime results are
based on an unoptimized implementation; we expect improved scalability with
mini-batch inference and mixed-precision training, but these optimizations must
be empirically validated.
In terms of broader impact, improved denoising can support biological discovery
and more stable downstream analyses, but over-imputation or misuse of denoised
matrices as ``ground truth'' can lead to spurious conclusions; we therefore
recommend reporting sensitivity to key hyperparameters and validating findings
against raw data and orthogonal assays when possible.

\section{Software and Data Availability}
For double-blind review, links to code and any processed artifacts are redacted.
An anonymized zip archive containing code, preprocessing scripts, and experiment
configurations is provided as a separate file with this submission for
reviewers. We will release the same materials publicly upon acceptance.
Synthetic datasets are generated with Splatter \citep{zappia2017splatter}.
The anonymized supplementary archive includes the code, configuration files,
dependency notes, and license information for released scripts and third-party
assets used to reproduce the reported experiments.

% Acknowledgments should only appear in the accepted (camera-ready) version.
% NeurIPS provides the {ack} environment (it is hidden automatically for
% submission). Keep it commented out for blind review.
% \begin{ack}
% \end{ack}



\clearpage
\bibliographystyle{plainnat}
\bibliography{references}

\clearpage
\appendix

\section{Synthetic Simulation Details}
\label{app:synthetic_details}
This appendix describes the data-generation protocol underlying
Table~\ref{tab:benchmark_core} and Appendix
Tables~\ref{tab:per_scenario_mse}--\ref{tab:benchmark_significance}. It
summarizes the fixed train/test/runtime splits, the 13 scenario definitions,
the grouped, batch-effect, and trajectory simulation settings, and the assays
and marker-gene annotations used by the benchmark metrics.
% Supplementary details for Section~\ref{sec:synthetic_data}.

\paragraph{Global settings and splits.}
All simulations start with $G=1000$ genes and are filtered to retain genes
expressed in at least three observed cells. We generate three splits:
(\textit{i}) a tuning split with $N=1000$ cells per scenario,
(\textit{ii}) a held-out test split with $N=5000$ cells per scenario, and
(\textit{iii}) a runtime-scaling split with
$N\in\{10000,20000,\ldots,100000\}$ on a single representative scenario
(\texttt{groups\_balanced\_moderate\_drop}). All datasets are generated with
fixed seeds (split-specific offsets), so each scenario corresponds to one
deterministic dataset per split.
We use the tuning split only for model selection and hyperparameter choices;
all reported reconstruction metrics are computed on the held-out test split.

\paragraph{Scenario catalog.}
The benchmark contains 13 scenarios: one no-dropout baseline
(\texttt{groups\_balanced\_nodrop}) and six scenario families, each with an
experiment-level dropout variant and a group-level dropout variant:
balanced groups (\texttt{groups\_balanced\_moderate\_drop}),
imbalanced groups (\texttt{groups\_imbalanced\_moderate\_drop}),
rare populations (\texttt{groups\_rare\_high\_drop}),
batch effects (\texttt{batch\_effects\_moderate\_drop}),
linear trajectories (\texttt{paths\_linear\_moderate\_drop}), and
branching trajectories (\texttt{paths\_branching\_moderate\_drop}).
Group-level variants append \texttt{\_group} to the scenario id.

\paragraph{Grouped population scenarios.}
Grouped datasets use Splat's \texttt{groups} mode with group proportions
(\texttt{group.prob}): balanced $(1/3,1/3,1/3)$, imbalanced $(0.6,0.3,0.1)$, and
rare $(0.50,0.25,0.20,0.05)$. We fix the fraction of differentially expressed
genes (\texttt{de.prob}) to $0.10$ in all grouped scenarios.
Dropout is simulated using Splatter's logistic dropout model, which maps log
mean expression to a dropout probability via a logistic function with midpoint
(\texttt{dropout.mid}) and slope (\texttt{dropout.shape}); we use
moderate-dropout settings $\texttt{dropout.mid}=3.0$ and
$\texttt{dropout.shape}=-1.0$, and rare high-dropout settings
$\texttt{dropout.mid}=4.0$ and $\texttt{dropout.shape}=-1.0$.
The scope of dropout heterogeneity is controlled by \texttt{dropout.type}:
experiment-level variants use \texttt{dropout.type=experiment} (shared
parameters), while group-level variants use \texttt{dropout.type=group} with
group-specific parameter vectors created by evenly perturbing the base values
across groups (\texttt{dropout.mid} range $\pm 0.4$, \texttt{dropout.shape}
range $\pm 0.25$).

\paragraph{Batch-effect scenario.}
The batch-effect family uses two batches of approximately equal size
(\texttt{batchCells}) and introduces batch-specific multiplicative gene factors
via the batch effect location/scale parameters
(\texttt{batch.facLoc}=0.10, \texttt{batch.facScale}=0.10). We also increase
library-size dispersion (\texttt{lib.scale}=0.40) to induce greater
heterogeneity in per-cell sequencing depth.

\paragraph{Trajectory scenarios.}
Trajectory datasets are generated with Splat's \texttt{paths} mode. Both linear
and branching trajectories use four groups with equal proportions
$(0.25,0.25,0.25,0.25)$, a higher fraction of DE genes (\texttt{de.prob}=0.20),
and stronger DE factor magnitudes via the log-normal location parameter
(\texttt{de.facLoc}=0.20). We set the number of interpolation steps
(\texttt{path.nSteps}) to 50. The trajectory topology is controlled by
\texttt{path.from}: the linear topology uses $\texttt{path.from}=(0,1,2,3)$,
while the branching topology uses $\texttt{path.from}=(0,1,1,3)$. Both
trajectory families use the same moderate dropout settings as the grouped
moderate-dropout scenarios, again with experiment-level
(\texttt{dropout.type=experiment}) and group-level (\texttt{dropout.type=group})
dropout variants.

\paragraph{Assays, normalization, and marker genes.}
Each dataset stores observed counts (\texttt{counts}) and pre-dropout ground
truth (\texttt{TrueCounts}), as well as CP10k-normalized log-transformed assays
$\texttt{logcounts}=\log_2(1+\mathrm{CP10k}(\texttt{counts}))$ and
$\texttt{logTrueCounts}=\log_2(1+\mathrm{CP10k}(\texttt{TrueCounts}))$, where
observed and true counts are normalized independently. Per-cell observed and
true library sizes and the normalization target sum (\texttt{targetSum}) are
stored in metadata. Marker genes are defined using \texttt{DEFacGroup*}
row-data columns: a gene is a marker if its differential-expression factor is
not equal to 1 in any group.



\clearpage
\section{Per-Scenario Synthetic Results}
\label{app:per_scenario}
This appendix expands the aggregate synthetic benchmark in
Table~\ref{tab:benchmark_core}. Table~\ref{tab:per_scenario_mse} reports
per-scenario overall MSE for each method, Table~\ref{tab:benchmark_marker}
summarizes marker-gene reconstruction errors, Figure~\ref{fig:mse_5000_bars}
visualizes the main synthetic metrics across methods,
Table~\ref{tab:benchmark_wins} reports scenario-level wins by metric, and
Table~\ref{tab:benchmark_significance} gives paired Wilcoxon tests comparing
MaskImpute with DCA across scenarios.
\begin{table}[H]
\caption{Per-scenario overall MSE on the synthetic test split. Entries report
mean $\pm$ repeat standard deviation over 5 seeds. Lower is better.}
\label{tab:per_scenario_mse}
\centering
\footnotesize
\resizebox{\textwidth}{!}{\begin{tabular}{@{}lrrrrrrrrr@{}}
\toprule
Scenario & MaskImpute & DCA & scVI & ALRA & MAGIC & AutoClass & ccImpute & SAVER & Baseline \\
\midrule
batch\_effects\_moderate\_drop & \textbf{0.2362 $\pm$ 0.0008} & 0.3505 $\pm$ 0.0005 & 0.3290 $\pm$ 0.0103 & 0.6831 $\pm$ 0.0000 & 0.6083 $\pm$ 0.0005 & 0.7074 $\pm$ 0.0078 & 0.6042 $\pm$ 0.0037 & 1.2058 $\pm$ 0.0002 & 1.5365 $\pm$ 0.0000 \\
batch\_effects\_moderate\_drop\_group & \textbf{0.2216 $\pm$ 0.0039} & 0.3243 $\pm$ 0.0005 & 0.3053 $\pm$ 0.0078 & 0.6628 $\pm$ 0.0000 & 0.5815 $\pm$ 0.0005 & 0.6794 $\pm$ 0.0040 & 0.5640 $\pm$ 0.0009 & 1.2284 $\pm$ 0.0002 & 1.5863 $\pm$ 0.0000 \\
groups\_balanced\_moderate\_drop & \textbf{0.2184 $\pm$ 0.0033} & 0.3428 $\pm$ 0.0004 & 0.3266 $\pm$ 0.0101 & 0.6444 $\pm$ 0.0000 & 0.6048 $\pm$ 0.0002 & 0.6844 $\pm$ 0.0059 & 0.6073 $\pm$ 0.0036 & 1.1731 $\pm$ 0.0002 & 1.4896 $\pm$ 0.0000 \\
groups\_balanced\_moderate\_drop\_group & \textbf{0.2104 $\pm$ 0.0023} & 0.3263 $\pm$ 0.0004 & 0.3148 $\pm$ 0.0087 & 0.6357 $\pm$ 0.0000 & 0.5750 $\pm$ 0.0002 & 0.6685 $\pm$ 0.0064 & 0.5516 $\pm$ 0.0006 & 1.1484 $\pm$ 0.0002 & 1.4636 $\pm$ 0.0000 \\
groups\_balanced\_nodrop & 0.1480 $\pm$ 0.0004 & 0.1488 $\pm$ 0.0004 & 0.1571 $\pm$ 0.0027 & 0.1470 $\pm$ 0.0000 & 0.1386 $\pm$ 0.0000 & 0.1518 $\pm$ 0.0012 & 0.0149 $\pm$ 0.0000 & 0.0208 $\pm$ 0.0000 & \textbf{0.0000 $\pm$ 0.0000} \\
groups\_imbalanced\_moderate\_drop & \textbf{0.2131 $\pm$ 0.0023} & 0.3270 $\pm$ 0.0004 & 0.3176 $\pm$ 0.0114 & 0.6413 $\pm$ 0.0000 & 0.6036 $\pm$ 0.0003 & 0.6938 $\pm$ 0.0054 & 0.5895 $\pm$ 0.0030 & 1.2456 $\pm$ 0.0001 & 1.6054 $\pm$ 0.0000 \\
groups\_imbalanced\_moderate\_drop\_group & \textbf{0.2092 $\pm$ 0.0014} & 0.2907 $\pm$ 0.0002 & 0.2824 $\pm$ 0.0063 & 0.5413 $\pm$ 0.0000 & 0.4840 $\pm$ 0.0004 & 0.5775 $\pm$ 0.0063 & 0.4191 $\pm$ 0.0008 & 0.9851 $\pm$ 0.0003 & 1.3147 $\pm$ 0.0000 \\
groups\_rare\_high\_drop & \textbf{0.3084 $\pm$ 0.0058} & 0.6620 $\pm$ 0.0010 & 0.6261 $\pm$ 0.0170 & 1.3608 $\pm$ 0.0000 & 1.3149 $\pm$ 0.0005 & 1.4241 $\pm$ 0.0053 & 1.8775 $\pm$ 0.0041 & 2.2323 $\pm$ 0.0003 & 2.5725 $\pm$ 0.0000 \\
groups\_rare\_high\_drop\_group & \textbf{0.2947 $\pm$ 0.0043} & 0.5961 $\pm$ 0.0007 & 0.5543 $\pm$ 0.0136 & 1.1689 $\pm$ 0.0000 & 1.1063 $\pm$ 0.0003 & 1.2197 $\pm$ 0.0036 & 1.5066 $\pm$ 0.0033 & 1.8991 $\pm$ 0.0003 & 2.2158 $\pm$ 0.0000 \\
paths\_branching\_moderate\_drop & \textbf{0.2197 $\pm$ 0.0020} & 0.3316 $\pm$ 0.0005 & 0.3184 $\pm$ 0.0029 & 0.6305 $\pm$ 0.0000 & 0.5930 $\pm$ 0.0004 & 0.6700 $\pm$ 0.0039 & 0.5823 $\pm$ 0.0033 & 1.2059 $\pm$ 0.0001 & 1.5366 $\pm$ 0.0000 \\
paths\_branching\_moderate\_drop\_group & \textbf{0.2153 $\pm$ 0.0023} & 0.3169 $\pm$ 0.0011 & 0.2961 $\pm$ 0.0049 & 0.6424 $\pm$ 0.0000 & 0.5757 $\pm$ 0.0002 & 0.6853 $\pm$ 0.0084 & 0.5893 $\pm$ 0.0072 & 1.2605 $\pm$ 0.0003 & 1.6292 $\pm$ 0.0000 \\
paths\_linear\_moderate\_drop & \textbf{0.2129 $\pm$ 0.0026} & 0.3311 $\pm$ 0.0003 & 0.3146 $\pm$ 0.0060 & 0.6210 $\pm$ 0.0000 & 0.5849 $\pm$ 0.0008 & 0.6621 $\pm$ 0.0066 & 0.5937 $\pm$ 0.0080 & 1.1616 $\pm$ 0.0001 & 1.4664 $\pm$ 0.0000 \\
paths\_linear\_moderate\_drop\_group & \textbf{0.2061 $\pm$ 0.0030} & 0.3289 $\pm$ 0.0008 & 0.3111 $\pm$ 0.0018 & 0.6131 $\pm$ 0.0000 & 0.5640 $\pm$ 0.0003 & 0.6419 $\pm$ 0.0033 & 0.5780 $\pm$ 0.0043 & 1.1043 $\pm$ 0.0002 & 1.3901 $\pm$ 0.0000 \\
\bottomrule
\end{tabular}
}
\end{table}

\begin{table}[H]
\caption{Marker-gene subset errors on the synthetic test split across 13
scenarios. Entries report means across scenario-level estimates; $\pm$ values
are standard deviations across scenarios. Lower is better.}
\label{tab:benchmark_marker}
\centering
\footnotesize
\begin{tabular}{@{}lrrr@{}}
\toprule
Method & Marker-MSE & Marker-MAE & Marker-gNRMSE \\
\midrule
MaskImpute & \textbf{0.2235 $\pm$ 0.0410} & \textbf{0.3470 $\pm$ 0.0331} & \textbf{1.1072 $\pm$ 0.0864} \\
DCA & 0.3603 $\pm$ 0.1296 & 0.4490 $\pm$ 0.0871 & 1.3265 $\pm$ 0.2055 \\
scVI & 0.3468 $\pm$ 0.1188 & 0.4364 $\pm$ 0.0775 & 1.2981 $\pm$ 0.1837 \\
ALRA & 0.6970 $\pm$ 0.2936 & 0.6471 $\pm$ 0.1560 & 1.7596 $\pm$ 0.3518 \\
MAGIC & 0.6440 $\pm$ 0.2867 & 0.6224 $\pm$ 0.1550 & 1.6879 $\pm$ 0.3479 \\
AutoClass & 0.6662 $\pm$ 0.0310 & 0.6383 $\pm$ 0.0245 & 1.7260 $\pm$ 0.0001 \\
ccImpute & 0.7454 $\pm$ 0.5446 & 0.5176 $\pm$ 0.2569 & 1.5869 $\pm$ 0.6903 \\
SAVER & 1.2286 $\pm$ 0.5043 & 0.6716 $\pm$ 0.2184 & 2.2907 $\pm$ 0.6536 \\
Baseline & 1.5151 $\pm$ 0.5781 & 0.7277 $\pm$ 0.2556 & 2.5166 $\pm$ 0.8063 \\
\bottomrule
\end{tabular}

\end{table}

\begin{figure}[H]
\centering
\includegraphics[width=\textwidth]{figures/mse_5000_bars.pdf}
\caption{Benchmark metrics on the synthetic test split ($N{=}5{,}000$ cells per
scenario). Each panel reports a metric aggregated across the 13 scenarios: bars
show the mean across scenarios and error bars show the standard deviation, with
lower error bars clipped at zero for non-negative metrics. Lower is better for
all panels.}
\label{fig:mse_5000_bars}
\end{figure}

\begin{table}[H]
\caption{Scenario-level wins on the synthetic test split. A win is assigned to
the method with the lowest scenario-level mean for the metric.}
\label{tab:benchmark_wins}
\centering
\footnotesize
\begin{tabular}{@{}ll@{}}
\toprule
Metric & Scenario-level wins \\
\midrule
MSE & MaskImpute: 12/13; Baseline: 1/13 \\
Dropout-MSE & MaskImpute: 12/12 \\
Biozero-MSE & Baseline: 13/13 \\
MAE & MaskImpute: 12/13; Baseline: 1/13 \\
gNRMSE & MaskImpute: 12/13; Baseline: 1/13 \\
Biozero-MAE & Baseline: 13/13 \\
Marker-MSE & MaskImpute: 12/13; Baseline: 1/13 \\
Marker-MAE & MaskImpute: 12/13; Baseline: 1/13 \\
Marker-gNRMSE & MaskImpute: 12/13; Baseline: 1/13 \\
\bottomrule
\end{tabular}

\end{table}

\begin{table}[H]
\caption{Paired Wilcoxon signed-rank tests comparing MaskImpute against DCA
across synthetic test scenarios. The one-sided test uses the alternative that
MaskImpute has lower error than DCA.}
\label{tab:benchmark_significance}
\centering
\footnotesize
\begin{tabular}{@{}lrr@{}}
\toprule
Metric & Wilcoxon two-sided $p$ & Wilcoxon MaskImpute $<$ DCA $p$ \\
\midrule
MSE & 0.0002 & 0.0001 \\
Dropout-MSE & 0.0005 & 0.0002 \\
Biozero-MSE & 0.0215 & 0.9915 \\
MAE & 0.0002 & 0.0001 \\
gNRMSE & 0.0005 & 0.0002 \\
Marker-MSE & 0.0002 & 0.0001 \\
Marker-MAE & 0.0002 & 0.0001 \\
Marker-gNRMSE & 0.0005 & 0.0002 \\
\bottomrule
\end{tabular}

\end{table}

\FloatBarrier
\clearpage
\section{Ablations, Calibration, and Synthetic Downstream Details}
\label{app:ablations}
This appendix reports diagnostic experiments for the MaskImpute design.
Table~\ref{tab:phase6_ablation} isolates the contribution of major components,
Table~\ref{tab:phase6_calibration} evaluates residual calibration strengths,
Table~\ref{tab:phase6_downstream} reports synthetic downstream clustering
diagnostics for selected outputs, and Table~\ref{tab:phase6_sensitivity}
summarizes one-scenario sensitivity to the shrinkage, biozero-weight, and zero
mask-rate settings.
\begin{table}[H]
\caption{MaskImpute component ablations on the synthetic test split. Entries are
means across 13 scenarios; $\pm$ values are 95\% confidence intervals across
scenarios for efficacy, MSE, and Biozero-MSE. ARI and NMI are
cluster-count-matched synthetic downstream scores. Lower is better for error
metrics; higher is better for ARI/NMI.}
\label{tab:phase6_ablation}
\centering
\footnotesize
\resizebox{\textwidth}{!}{\input{generated/phase6_ablation_table.tex}}
\end{table}

\begin{table}[H]
\caption{Label-free residual calibration on the synthetic test split. The
calibrated output is $(1-\lambda)X_{\mathrm{MI}}+\lambda X_{\mathrm{obs}}$.
Entries are means across 13 scenarios. Lower is better for error metrics;
higher is better for ARI/NMI.}
\label{tab:phase6_calibration}
\centering
\footnotesize
\resizebox{\textwidth}{!}{\begin{tabular}{lrrrrrr}
\toprule
Output & Efficacy & MSE & Biozero-MSE & Dropout-MSE & ARI & NMI \\
\midrule
MaskImpute & 0.251 $\pm$ 0.017 & 0.224 & 0.013 & 0.304 & 0.718 & 0.764 \\
$0.95X_{MI}+0.05X_{obs}$ & 0.243 $\pm$ 0.018 & 0.219 & 0.012 & 0.310 & 0.723 & 0.766 \\
$0.90X_{MI}+0.10X_{obs}$ & 0.243 $\pm$ 0.021 & 0.221 & 0.011 & 0.331 & 0.740 & 0.772 \\
$0.85X_{MI}+0.15X_{obs}$ & 0.250 $\pm$ 0.025 & 0.230 & 0.010 & 0.368 & 0.721 & 0.760 \\
$0.80X_{MI}+0.20X_{obs}$ & 0.264 $\pm$ 0.030 & 0.247 & 0.009 & 0.420 & 0.730 & 0.769 \\
$0.75X_{MI}+0.25X_{obs}$ & 0.287 $\pm$ 0.037 & 0.271 & 0.008 & 0.488 & 0.725 & 0.767 \\
Observed input & 1.524 $\pm$ 0.314 & 1.524 & 0.000 & 3.366 & 0.642 & 0.663 \\
\bottomrule
\end{tabular}
}
\end{table}

\begin{table}[H]
\caption{Synthetic downstream clustering diagnostic for selected outputs.
Matched columns select the clustering grid point with the closest number of
clusters to the simulator-provided label count; grid columns average across the
full PCA--SNN--Leiden grid. Higher is better.}
\label{tab:phase6_downstream}
\centering
\footnotesize
\input{generated/phase6_downstream_table.tex}
\end{table}

\begin{table}[H]
\caption{One-scenario sensitivity analysis on
\texttt{groups\_balanced\_moderate\_drop}. Efficacy is
MSE $+2\times$ Biozero-MSE. Lower is better for error metrics; higher is better
for ARI.}
\label{tab:phase6_sensitivity}
\centering
\footnotesize
\resizebox{\textwidth}{!}{\begin{tabular}{lrrrrrr}
\toprule
Parameter & Value & Efficacy & MSE & Biozero-MSE & Dropout-MSE & ARI \\
\midrule
Shrinkage $\alpha$ & 0 & 0.255 & 0.221 & 0.017 & 0.292 & 0.996 \\
Shrinkage $\alpha$ & 0.1 & 0.250 & 0.218 & 0.016 & 0.287 & 0.995 \\
Shrinkage $\alpha$ & 0.2 & 0.246 & 0.217 & 0.014 & 0.285 & 0.995 \\
Shrinkage $\alpha$ & 0.3 & 0.245 & 0.218 & 0.013 & 0.287 & 0.995 \\
Shrinkage $\alpha$ & 0.5 & 0.248 & 0.227 & 0.011 & 0.306 & 0.993 \\
Shrinkage $\alpha$ & 0.7 & 0.264 & 0.247 & 0.008 & 0.350 & 0.993 \\
Shrinkage $\alpha$ & 1 & 0.319 & 0.308 & 0.005 & 0.479 & 0.992 \\
Biozero weight $\gamma$ & 0 & 0.361 & 0.166 & 0.097 & 0.184 & 0.989 \\
Biozero weight $\gamma$ & 0.5 & 0.241 & 0.192 & 0.025 & 0.244 & 0.998 \\
Biozero weight $\gamma$ & 1 & 0.242 & 0.211 & 0.015 & 0.276 & 0.996 \\
Biozero weight $\gamma$ & 1.2 & 0.245 & 0.218 & 0.013 & 0.287 & 0.995 \\
Biozero weight $\gamma$ & 2 & 0.260 & 0.243 & 0.008 & 0.328 & 0.995 \\
Biozero weight $\gamma$ & 5 & 0.307 & 0.300 & 0.004 & 0.414 & 0.983 \\
Zero mask rate $p_0$ & 0.005 & 0.244 & 0.216 & 0.014 & 0.284 & 0.996 \\
Zero mask rate $p_0$ & 0.01 & 0.245 & 0.218 & 0.013 & 0.287 & 0.995 \\
Zero mask rate $p_0$ & 0.02 & 0.246 & 0.222 & 0.012 & 0.293 & 0.995 \\
Zero mask rate $p_0$ & 0.05 & 0.251 & 0.229 & 0.011 & 0.307 & 0.991 \\
\bottomrule
\end{tabular}
}
\end{table}

\FloatBarrier
\clearpage
\section{Real-Data Label Agreement Details}
\label{app:real_data}
This appendix contains the real-data label-agreement checks used as downstream
diagnostics rather than direct imputation-accuracy evidence.
Table~\ref{tab:real_data} reports grid-averaged and cluster-count-matched ARI
and NMI, together with batch ARI, for Baron, Zeisel, and PBMC68k after each
imputation method.
\begin{table}[H]
\caption{Real-data label agreement after imputation. ARI grid and NMI grid are
means over five seeds and a PCA--SNN--Leiden robustness grid. Matched columns
select the grid point whose number of clusters is closest to the number of
annotated cell types for each seed. Higher is better for ARI and NMI. Batch ARI
measures cluster--batch association, where lower is better.}
\label{tab:real_data}
\centering
\footnotesize
\resizebox{\textwidth}{!}{\begin{tabular}{@{}llccccc@{}}
\toprule
Dataset & Method & ARI grid & NMI grid & ARI matched & NMI matched & Batch ARI \\
\midrule
Baron & Observed & 0.508 $\pm$ 0.170 & 0.771 $\pm$ 0.059 & 0.667 $\pm$ 0.098 & 0.830 $\pm$ 0.029 & 0.147 $\pm$ 0.040 \\
Baron & DCA & 0.442 $\pm$ 0.158 & 0.746 $\pm$ 0.062 & 0.607 $\pm$ 0.031 & 0.818 $\pm$ 0.005 & 0.160 $\pm$ 0.038 \\
Baron & scVI & 0.424 $\pm$ 0.139 & 0.739 $\pm$ 0.059 & 0.607 $\pm$ 0.031 & 0.816 $\pm$ 0.010 & 0.162 $\pm$ 0.037 \\
Baron & ALRA & 0.390 $\pm$ 0.155 & 0.721 $\pm$ 0.061 & 0.619 $\pm$ 0.031 & 0.811 $\pm$ 0.005 & 0.116 $\pm$ 0.033 \\
Baron & MAGIC & 0.315 $\pm$ 0.116 & 0.688 $\pm$ 0.053 & 0.523 $\pm$ 0.000 & 0.779 $\pm$ 0.000 & 0.156 $\pm$ 0.051 \\
Baron & MaskImpute & 0.456 $\pm$ 0.146 & 0.754 $\pm$ 0.059 & 0.684 $\pm$ 0.026 & 0.845 $\pm$ 0.007 & 0.163 $\pm$ 0.030 \\
\addlinespace
Zeisel & Observed & 0.409 $\pm$ 0.174 & 0.629 $\pm$ 0.045 & 0.757 $\pm$ 0.002 & 0.738 $\pm$ 0.002 & 0.038 $\pm$ 0.005 \\
Zeisel & DCA & 0.381 $\pm$ 0.156 & 0.611 $\pm$ 0.035 & 0.655 $\pm$ 0.023 & 0.670 $\pm$ 0.012 & 0.042 $\pm$ 0.008 \\
Zeisel & scVI & 0.374 $\pm$ 0.152 & 0.604 $\pm$ 0.036 & 0.650 $\pm$ 0.014 & 0.675 $\pm$ 0.005 & 0.042 $\pm$ 0.007 \\
Zeisel & ALRA & 0.327 $\pm$ 0.126 & 0.610 $\pm$ 0.032 & 0.601 $\pm$ 0.015 & 0.676 $\pm$ 0.004 & 0.040 $\pm$ 0.005 \\
Zeisel & MAGIC & 0.248 $\pm$ 0.135 & 0.587 $\pm$ 0.040 & 0.623 $\pm$ 0.027 & 0.692 $\pm$ 0.009 & 0.055 $\pm$ 0.010 \\
Zeisel & MaskImpute & 0.366 $\pm$ 0.169 & 0.621 $\pm$ 0.042 & 0.688 $\pm$ 0.007 & 0.700 $\pm$ 0.009 & 0.042 $\pm$ 0.008 \\
\addlinespace
PBMC68k & Observed & 0.214 $\pm$ 0.030 & 0.424 $\pm$ 0.049 & 0.216 $\pm$ 0.002 & 0.452 $\pm$ 0.022 & 0.000 $\pm$ 0.000 \\
PBMC68k & DCA & 0.212 $\pm$ 0.061 & 0.430 $\pm$ 0.052 & 0.277 $\pm$ 0.023 & 0.486 $\pm$ 0.017 & 0.001 $\pm$ 0.000 \\
PBMC68k & scVI & 0.216 $\pm$ 0.048 & 0.424 $\pm$ 0.049 & 0.254 $\pm$ 0.022 & 0.460 $\pm$ 0.009 & 0.000 $\pm$ 0.000 \\
PBMC68k & ALRA & 0.169 $\pm$ 0.054 & 0.412 $\pm$ 0.026 & 0.281 $\pm$ 0.032 & 0.462 $\pm$ 0.004 & 0.000 $\pm$ 0.000 \\
PBMC68k & MAGIC & 0.145 $\pm$ 0.045 & 0.374 $\pm$ 0.027 & 0.206 $\pm$ 0.001 & 0.449 $\pm$ 0.004 & 0.000 $\pm$ 0.000 \\
PBMC68k & MaskImpute & 0.187 $\pm$ 0.038 & 0.402 $\pm$ 0.030 & 0.220 $\pm$ 0.027 & 0.419 $\pm$ 0.009 & 0.000 $\pm$ 0.000 \\
\bottomrule
\end{tabular}
}
\end{table}

\FloatBarrier
\clearpage
\section{Runtime Scaling Details}
\label{app:runtime}
Runtime artifacts were generated on an NVIDIA DGX H100 system, with all
methods in the runtime comparison benchmarked on the same hardware. Full CPU,
RAM, and software-version details should be finalized with the post-feedback
benchmark rerun.
Figure~\ref{fig:runtime_linear} plots wall-clock scaling across synthetic
cell-count settings, and Table~\ref{tab:runtime_scaling} provides the
corresponding runtime values in seconds for each method and cell count.
\begin{figure}[H]
\centering
\includegraphics[width=\textwidth]{figures/runtime_scaling_linear.png}
\caption{Runtime scaling on the synthetic scaling split (linear scale). Each
point reports wall-clock runtime for a single dataset at the corresponding cell
count.}
\label{fig:runtime_linear}
\end{figure}

\begin{table}[H]
\caption{Runtime (seconds) for the synthetic scaling split.}
\label{tab:runtime_scaling}
\centering
\footnotesize
\resizebox{\textwidth}{!}{\begin{tabular}{@{}lrrrrrrrrrr@{}}
\toprule
Method & 10000 & 20000 & 30000 & 40000 & 50000 & 60000 & 70000 & 80000 & 90000 & 100000 \\
\midrule
MaskImpute & 295.15 & 559.69 & 830.76 & 1134.28 & 1414.03 & 1692.52 & 1909.26 & 2165.32 & 2430.20 & 2838.32 \\
DCA & 208.71 & 476.29 & 570.12 & 651.25 & 988.29 & 1037.25 & 1330.96 & 1549.41 & 1399.46 & 1606.54 \\
MAGIC & 2.77 & 4.98 & 7.15 & 10.61 & 15.04 & 20.31 & 25.26 & 31.02 & 36.74 & 43.53 \\
AutoClass & 175.62 & 298.51 & 577.81 & 819.63 & 1281.41 & 2727.90 & 2142.30 & 2212.01 & 2719.49 & 3438.38 \\
ccImpute & 52.90 & 126.05 & 215.93 & 337.61 & 522.05 & 691.23 & 966.11 & 1172.14 & 1647.80 & 1811.54 \\
SAVER & 139.83 & 346.51 & 417.33 & 718.73 & 790.14 & 1007.26 & 1204.63 & 1287.87 & 1435.66 & 2132.14 \\
Baseline & 11.10 & 23.73 & 37.34 & 52.66 & 64.58 & 77.05 & 91.47 & 104.35 & 116.88 & 132.26 \\
\bottomrule
\end{tabular}
}
\end{table}

\FloatBarrier
\clearpage
\input{checklist.tex}

\end{document}
